\chapter{Saving of Results} 

\section{Node Results}
By default, the degree of freedom is saved to the universal file. But
sometimes one needs (additionally) the time derivative (e.g. in
\analysis{acoustics}, to calculate the pressure from the acoustic velocity
potential). Herefore, the command is given for \analysis{mechanics}

\begin{verbatim}
<storeResults>
   <nodeResult type="mechVelocity">
    <allRegions/>
   </nodeResult>
</storeResults>
\end{verbatim}

may be used. The value of the attribute \attribute{type} is PDE dependent,
e.g. for mechanics, the following values are allowed:
\begin{itemize}
\item \attrivalue{mechDisplacement}
\item \attrivalue{mechVelocity}
\item \attrivalue{mechAcceleration}
\end{itemize}
The allowed types are mentioned at the individual PDE sections.

By default, the values on every node are stored. With the additional
attribute \attribute{region} one can restrict the saved results to one or
more selected regions.
To store the results of the regions \attrivalue{bar} and \attrivalue{gap}, one has to write
\begin{verbatim}
<storeResults>
   <nodeResult type="mechVelocity">
    <regionList>
     <region name="bar"/>
     <region name="gap"/>
    </regionList>
   </nodeResult>
</storeResults>
\end{verbatim}


\section{Element Results}
In general, saving of element results is, beside the different keyword
{\bf elemResults} totally similar to the procedure of saving nodal
results:
\begin{verbatim}
<storeResults>
   <elemResult type="mechStress"/>
</storeResults>
\end{verbatim}
Again, the values of \attribute{type} depend on the regarded PDE. 


\section{Results at Special Points} 
In many applications we are only interested in the response of the
system at a single point. In order to reduce necessary disk space,
it is possible to store the nodal results only at a discrete
location. Therefore, the command \element{nodeList} with the
\subelement{nodes}, with attribute \attribute{name} and attribute value as node name \attrivalue{nodeName} is provided. This command
reduces the amount of results stored tremendously, because no
universal-file is written.

In the config-file, the command is as follows

\begin{verbatim}
<storeResults>
   <nodeResult type="mechVelocity">
    <nodeList>
     <nodes name="nodeName"/>
    </nodeList>
   </nodeResult>
</storeResults>
\end{verbatim}

has to be used. The results are stored in the -- if necessary new
created -- directory {\em history} in a file with the name {\em
  'xml-file-name'-'quantity'-'nodeNr'}.hist.

\begin{verbatim}
nsel,s,loc,x,len-1e-6,len+1e-6
nsel,r,loc,y,height
wsavnod,'ptName'
\end{verbatim}

As a special case, for storing nodal results at specific radial distances
from a center point (e.g for directivity plots), first the nodes at these
points have to be defined (searched) using the element \element{directivityNodes} in section \topelement{$<$domain$>$}. Then the
procedure above can be used, using the name given to this set of
nodes. 

\section{Results at Special Elements} 
Again, saving results of selected elements is similar to the procedure
of saving results on selected nodes. The commands name is \element{elemList}, and the necessary keyword attribute name \attribute{name} with value \attrivalue{elName}:
\begin{verbatim}
   <storeResults>
      <elemResult type="mechStress">
       <elemList>
        <elem name="elName"/>
       </elemList>
      </elemResult>
   </storeResults>
\end{verbatim}

